\section{Regularity of the limit densities}\label{app:regularity}

We prove the regularity used in Sections~\ref{sec:bulk-applications}
and~\ref{app:modes}.

\begin{lemma}[entire densities and parameter continuity]\label{lem:analytic-density}
For every $a\in(\frac12,1)$ and $1<r<1/a$, the characteristic function
$\phi_a$ of $\mathcal L_a$ satisfies
\[
 |\phi_a(t)|\le 2e^{-c|t|^r}\qquad(t\in\mathbb R)
\]
for some $c=c(a,r)>0$. The same $c$ works on a sufficiently small
neighborhood of any fixed $a$. For every compact
$J\subset(1/2,1)$ and every $\beta>0$, there is $C_{J,\beta}<\infty$ such that
\[
 |\partial_t\phi_a(t)|+|\partial_a\phi_a(t)|
 \le C_{J,\beta}(1+|t|)^{-\beta}
 \qquad(a\in J,\ t\in\mathbb R).
\]
For every integer $k\ge0$, the derivatives $\partial_u^k h_a(u)$ and
$\partial_a\partial_u^k h_a(u)$ exist and are jointly continuous in $(a,u)$.
For every $q\in[0,1]$, the density of the limit with first-step bias $q$
extends to an entire function and is not constant on an interval of
positive length.
\end{lemma}

\begin{proof}
We first check weak continuity in $a$ to obtain a Fourier bound local in memory. On a compact
$J\subset(1/2,1)$, the martingale
$\mathcal M_n(a)=\Gamma(n)S_n/\Gamma(n+a)$ of~\cite[(2.4), Appendix~A]{bercu2018} satisfies
\[
 \sup_{a\in J}\E|\mathcal L_a-\mathcal M_n(a)|^2
 \le C_J\sum_{k\ge n}k^{-2\min J}\longrightarrow0.
\]
The walk's centered next step has conditional variance at most one,
and the gamma ratios are uniformly bounded on $J$. The finite laws are
continuous in $a$, so $\phi_a(t)$ is jointly continuous, uniformly on
compact frequency intervals.

The fixed-point equation of~\cite[Theorem~1.1 and (2.10)]{glr2026} gives
\[
 |\phi_a(t)|\le\int_0^1|\phi_a(tv^a)|\,
                  |\phi_a(t(1-v)^a)|\,dv.
\]
Here $\phi_a$ is the characteristic function conditional on a positive
first step. The random sign disappears after taking absolute values, since
$|\phi_a(-s)|=|\phi_a(s)|$.
Fix $a_0$ and $1<r<1/a_0$, and choose $J$ about $a_0$ with $r\max J<1$.
Put $D(v)=v^{r\max J}+(1-v)^{r\max J}-1$. We have $D(v)>0$ for
$0<v<1$. Dominated convergence therefore permits $K>1$ such that
\[
 R:=2\int_0^1e^{-K^r(\log2)D(v)}\,dv<1.
\]
The law of $\mathcal L_{a_0}$ has a density~\cite[Theorem~1.3]{glr2026}, so
the Riemann--Lebesgue lemma gives $T>0$ with
$|\phi_{a_0}(t)|<2^{-K^r}$ for $|t|\ge T$. Joint continuity, after
shrinking $J$, gives $|\phi_a(t)|\le2^{1-K^r}$ for $a\in J$ and
$T\le|t|\le KT$. Set
$c=(\log2)/T^r$ and $b(t)=2e^{-ct^r}$ for $t\ge0$.
On $[0,T]$ we have $b\ge1$, and on $[T,KT]$ we have
$b\ge2^{1-K^r}$. Thus $|\phi_a|\le b$ on $[0,KT]$.

For $t\ge KT$,
\[
 \frac{\int_0^1b(tv^a)b(t(1-v)^a)\,dv}{b(t)}
 \le2\int_0^1e^{-ct^rD(v)}\,dv\le R<1.
\]
Suppose the bound $|\phi_a|\le b$ stops at a finite
$t_*\ge KT$. By continuity it still holds on $[0,t_*]$.
But at $t_*$ the fixed-point inequality gives
$|\phi_a(t_*)|\le Rb(t_*)<b(t_*)$, which is strict. By continuity the
bound holds a little past $t_*$, a contradiction. Hence it holds on
$[0,\infty)$, and on $\mathbb R$ by evenness.

Since $r>1$, the Fourier integral converges absolutely and locally
uniformly for complex arguments. Indeed, for every $B<\infty$,
$\int_{\mathbb R}e^{B|t|}|\phi_a(t)|\,dt<\infty$.
It therefore defines an entire function.
After any fixed number of spatial derivatives, the bound on $J$
still supplies an integrable majorant, giving joint continuity by
dominated convergence. For initial bias $q$, the triangle inequality
applied to $q\phi_a(t)+(1-q)\phi_a(-t)$ gives the same bound.
A density of integral one cannot be constant on $\mathbb R$.
If it were constant on an interval, the identity theorem would
make it constant everywhere.

To differentiate in memory, we first work on bounded frequency intervals.
The moment recursion of~\cite[Theorem~1.4]{glr2026} is
\[
 m_0=m_1=1,\qquad
 m_n(z)=\frac{\sum_{j=1}^{n-1}\theta_j(z)m_j(z)m_{n-j}(z)}
                   {nz-\theta_n(z)},\qquad
 \theta_j(z)=\begin{cases}1&j\text{ even},\\z&j\text{ odd}.
                  \end{cases}
\]
Choose a closed complex disc about $a_0$ on which
$\operatorname{Re}z\ge\alpha>1/2$, $|z|<1$ and
$|\operatorname{Im}z|\le\rho$. All denominators for $n\ge2$ have
real part at least $n\alpha-1>0$.
Put $A=2\alpha-1$ and $B=A^{-1}$. Induction gives
\[
 |m_n(z)|\le AB^n\quad(n\ge1),\qquad
 \frac{(n-1)A^2}{n\alpha-1}\le A\quad(n\ge2).
\]
For $x>0$ and real $y$, the gamma product gives
\[
 \left|\frac{\Gamma(x)}{\Gamma(x+iy)}\right|^2
 =\prod_{j\ge0}\left(1+\frac{y^2}{(x+j)^2}\right)
 \le\exp\left(\frac{y^2}{x^2}+\frac{y^2}{x}\right).
\]
Taking $x=1+n\operatorname{Re}z$ and $y=n\operatorname{Im}z$ gives
\[
 \left|\frac{m_n(z)}{\Gamma(1+nz)}\right|
 \le\frac{C B_0^n}{\Gamma(1+n\alpha)}.
\]
Here $C,B_0$ are independent of $n,z$. We used monotonicity of
$\Gamma$ on $[2,\infty)$ for $n\ge2$, and enlarged $C$ for $n=0,1$.
Thus $\sum_{n\ge0}m_n(z)(it)^n/\Gamma(1+nz)$ converges locally
uniformly in complex $t$ and $z$ in the disc's interior.
For real $z=a$ this is $\phi_a(t)$, by the moment formula and
exponential moments of~\cite[Corollary~2.11]{glr2026}.
Cauchy's formula gives local joint continuity of every frequency
derivative and its memory derivative.

We now obtain bounds at large frequencies without a numerical input.
Shrink $J$ about $a_0$ so that it lies inside this disc and retains the
uniform exponential bound. For $t\ge0$ write
$f_a(t)=|\phi_a(t)|$, $u_a(t)=|\partial_t\phi_a(t)|$ and
$w_a(t)=|\partial_a\phi_a(t)|$. The smoothing equation is
\[
 \phi_a(t)=\int_0^1\phi_a(tv^a)
 \{p\phi_a(t(1-v)^a)+(1-p)\phi_a(-t(1-v)^a)\}\,dv,
 \qquad p=(1+a)/2.
\]
The local bounds just proved justify differentiation under this integral,
since $v^{\min J}|\log v|$ is integrable. Put $x=tv^a$,
$y=t(1-v)^a$ and, for nonnegative $z$, define
\[
 (\mathcal K_a z)(t)=\int_0^1\{z(x)f_a(y)+f_a(x)z(y)\}\,dv.
\]
Taking absolute values and discarding factors $v^a,(1-v)^a\le1$
in the frequency derivative gives
\[
 u_a\le\mathcal K_a u_a,\qquad
 w_a\le\mathcal K_a w_a+S_a,
\]
where
\[
 S_a(t)=\int_0^1\{x|\log v|u_a(x)f_a(y)
       +y|\log(1-v)|f_a(x)u_a(y)+f_a(x)f_a(y)\}\,dv.
\]
The last term bounds the derivative of $p$, since $p'=1/2$ and
$|\phi_a(y)-\phi_a(-y)|\le2f_a(y)$.

Fix $\beta>0$. Splitting at $1/2$ and reflecting the second half gives
\[
 K_{a,\beta}(t)=2\int_0^{1/2}
 \left\{\left(\frac{1+t}{1+x}\right)^\beta f_a(y)
       +\left(\frac{1+t}{1+y}\right)^\beta f_a(x)\right\}\,dv
 \longrightarrow0
\]
uniformly for $a\in J$ as $t\to\infty$.
Indeed, $y\ge t/2$ bounds the first term by
$2(1+t)^\beta e^{-c(t/2)^r}$ and the second ratio by $2^\beta$.
Also $x\ge tv^{\max J}$, so the exponential bound and dominated
convergence make the second integral tend to zero uniformly.
Choose $T$ so that $K_{a,\beta}(t)\le1/2$ for all $a\in J$, $t\ge T$.
For $X\ge T$ let
\[
 U_X=\max_{a\in J,\,0\le t\le X}(1+t)^\beta u_a(t).
\]
All arguments $x,y$ are at most $t$. The first smoothing inequality gives
\[
 U_X\le\max\{U_T,U_X/2\},\qquad\text{hence }U_X\le U_T.
\]
The bound is independent of $X$. This proves arbitrary polynomial decay for $u_a$.

The same decay holds for $S_a$. On $0<v\le1/2$, the first
logarithmic term uses the exponential decay of $f_a(y)$ and boundedness
of $u_a(x)$, with $x|\log v|\le tv^{\min J}|\log v|$.
The second uses arbitrary polynomial decay of $u_a(y)$,
$f_a(x)\le1$, and $|\log(1-v)|\le\log2$. The last uses the
exponential decay of $f_a(y)$. Reflection covers the other half.
Thus $S_\beta:=\sup_{a\in J,t\ge0}(1+t)^\beta S_a(t)<\infty$.
Defining $W_X$ as above with $w_a$ in place of $u_a$, we obtain
\[
 W_X\le\max\{W_T,W_X/2+S_\beta\}
 \quad\Longrightarrow\quad W_X\le\max\{W_T,2S_\beta\}.
\]
Conjugate symmetry extends both derivative bounds to negative $t$.
A finite cover gives the stated bounds on every compact memory interval.
Since $\beta$ is arbitrary, multiplication by any fixed power of $t$
leaves an integrable uniform majorant. Differentiation under Fourier
inversion and dominated convergence prove the claimed joint continuity
of $\partial_a\partial_u^k h_a$ for every $k$.
\end{proof}
